\documentclass[reqno]{amsart}
\usepackage[T1]{fontenc}
\usepackage{amsmath,amssymb,amsthm}
\usepackage{xr-hyper}
\usepackage[hidelinks]{hyperref}
\externaldocument[ob-]{20261126T000000--the-objects-of-the-construction__workbook}
\externaldocument[cm-]{20261124T223100--complex-multiplication-and-inert-reduction__workbook}

% Chapter number; \ref{chap} gives it.
\newcounter{chap}
% Chapter numbers, generated from 20261125T012000--grh-workbook__guide.org
% by scripts/make-chapters.py.  Do not edit: change the order of the
% index and rerun the script (./preview-site.sh and CI do this).
%
% Usage in a sheet, after \newcounter{chap}:
%   % Chapter numbers, generated from 20261125T012000--grh-workbook__guide.org
% by scripts/make-chapters.py.  Do not edit: change the order of the
% index and rerun the script (./preview-site.sh and CI do this).
%
% Usage in a sheet, after \newcounter{chap}:
%   % Chapter numbers, generated from 20261125T012000--grh-workbook__guide.org
% by scripts/make-chapters.py.  Do not edit: change the order of the
% index and rerun the script (./preview-site.sh and CI do this).
%
% Usage in a sheet, after \newcounter{chap}:
%   \input{chapters}\chapnumber{<sheet ID>}
\makeatletter
\providecommand\chapters@set[2]{\@namedef{chapters@#1}{#2}}
\chapters@set{20261125T001000}{1}% zeros-to-ramp-sums
\chapters@set{20261125T001100}{2}% inert-primes-and-estimate-b
\chapters@set{20261125T005000}{3}% galois-fourier-and-averaging
\chapters@set{20261125T005500}{4}% columns-from-a-sphere
\chapters@set{20261125T000000}{5}% slope-method
\chapters@set{20261126T030100}{6}% a-twist-for-each-prime
\chapters@set{20261125T003000}{7}% the-chain-and-the-square-root
\chapters@set{20261125T010000}{8}% rank-lemma-and-zero-estimate
\chapters@set{20261125T010500}{9}% poincare-series-and-near-orthogonality
\chapters@set{20261124T223000}{10}% elliptic-curves-for-the-proof
\chapters@set{20261125T009500}{11}% fock-model-and-theta-intertwiners
\chapters@set{20261124T223100}{12}% complex-multiplication-and-inert-reduction
\chapters@set{20261125T008000}{13}% one-cm-determinant
\chapters@set{20261124T214500}{14}% supersingular-theta-divisibility
\chapters@set{20261124T214600}{15}% the-isotropic-plane
\chapters@set{20261125T013700}{16}% pairings-on-a-formal-group-of-height-two
\chapters@set{20261125T013500}{17}% hecke-nondegeneracy-on-the-tree
\chapters@set{20261126T030200}{18}% the-bound-is-attained
\chapters@set{20261125T009000}{19}% kronecker-limit-and-chowla-selberg
\chapters@set{20261125T002000}{20}% definite-theta-lifts
\chapters@set{20261125T015500}{21}% siegel-weil-in-the-smallest-case
\chapters@set{20261125T002100}{22}% definite-theta-lifts-and-norms
\chapters@set{20261125T016500}{23}% hecke-operators-and-theta-lifts
\chapters@set{20261125T016700}{24}% divisibility-along-the-chain
\chapters@set{20261125T004000}{25}% denominators-and-sha
\chapters@set{20261125T004100}{26}% tate-twists-and-sha
\chapters@set{20261126T020430}{27}% the-minus-part-by-gauss-sums
\chapters@set{20261125T014100}{28}% congruence-modules
\chapters@set{20261125T014500}{29}% eisenstein-congruences-and-selmer-classes
\chapters@set{20261125T014700}{30}% pseudocharacters-and-the-length-bound
\chapters@set{20261125T011000}{31}% galois-representations-black-box
\chapters@set{20261126T030300}{32}% the-data
\chapters@set{20261126T020210}{33}% bessel-poincare-series-and-their-gram-operators
\chapters@set{20261126T030400}{34}% the-seesaw-identity
\chapters@set{20261126T030500}{35}% divisibility-of-the-composite-lift
\chapters@set{20261126T030600}{36}% theorem-a-short-forms
\chapters@set{20261126T030700}{37}% the-determinant-of-a-chain
\chapters@set{20261126T030800}{38}% the-cost-of-one-lift
\chapters@set{20261126T030900}{39}% denominators-of-projections-to-lifts
\chapters@set{20261126T031000}{40}% twists-with-few-congruences
\chapters@set{20261126T031100}{41}% theorem-c-and-the-main-inequality
\providecommand\chapnumber[1]{%
  \@ifundefined{chapters@#1}%
    {\PackageError{chapters}{Sheet #1 is not listed in the index}%
      {List it in 20261125T012000--grh-workbook__guide.org\MessageBreak
       and run scripts/make-chapters.py.}}%
    {\edef\thechap{\@nameuse{chapters@#1}}}}
\makeatother
\chapnumber{<sheet ID>}
\makeatletter
\providecommand\chapters@set[2]{\@namedef{chapters@#1}{#2}}
\chapters@set{20261125T001000}{1}% zeros-to-ramp-sums
\chapters@set{20261125T001100}{2}% inert-primes-and-estimate-b
\chapters@set{20261125T005000}{3}% galois-fourier-and-averaging
\chapters@set{20261125T005500}{4}% columns-from-a-sphere
\chapters@set{20261125T000000}{5}% slope-method
\chapters@set{20261126T030100}{6}% a-twist-for-each-prime
\chapters@set{20261125T003000}{7}% the-chain-and-the-square-root
\chapters@set{20261125T010000}{8}% rank-lemma-and-zero-estimate
\chapters@set{20261125T010500}{9}% poincare-series-and-near-orthogonality
\chapters@set{20261124T223000}{10}% elliptic-curves-for-the-proof
\chapters@set{20261125T009500}{11}% fock-model-and-theta-intertwiners
\chapters@set{20261124T223100}{12}% complex-multiplication-and-inert-reduction
\chapters@set{20261125T008000}{13}% one-cm-determinant
\chapters@set{20261124T214500}{14}% supersingular-theta-divisibility
\chapters@set{20261124T214600}{15}% the-isotropic-plane
\chapters@set{20261125T013700}{16}% pairings-on-a-formal-group-of-height-two
\chapters@set{20261125T013500}{17}% hecke-nondegeneracy-on-the-tree
\chapters@set{20261126T030200}{18}% the-bound-is-attained
\chapters@set{20261125T009000}{19}% kronecker-limit-and-chowla-selberg
\chapters@set{20261125T002000}{20}% definite-theta-lifts
\chapters@set{20261125T015500}{21}% siegel-weil-in-the-smallest-case
\chapters@set{20261125T002100}{22}% definite-theta-lifts-and-norms
\chapters@set{20261125T016500}{23}% hecke-operators-and-theta-lifts
\chapters@set{20261125T016700}{24}% divisibility-along-the-chain
\chapters@set{20261125T004000}{25}% denominators-and-sha
\chapters@set{20261125T004100}{26}% tate-twists-and-sha
\chapters@set{20261126T020430}{27}% the-minus-part-by-gauss-sums
\chapters@set{20261125T014100}{28}% congruence-modules
\chapters@set{20261125T014500}{29}% eisenstein-congruences-and-selmer-classes
\chapters@set{20261125T014700}{30}% pseudocharacters-and-the-length-bound
\chapters@set{20261125T011000}{31}% galois-representations-black-box
\chapters@set{20261126T030300}{32}% the-data
\chapters@set{20261126T020210}{33}% bessel-poincare-series-and-their-gram-operators
\chapters@set{20261126T030400}{34}% the-seesaw-identity
\chapters@set{20261126T030500}{35}% divisibility-of-the-composite-lift
\chapters@set{20261126T030600}{36}% theorem-a-short-forms
\chapters@set{20261126T030700}{37}% the-determinant-of-a-chain
\chapters@set{20261126T030800}{38}% the-cost-of-one-lift
\chapters@set{20261126T030900}{39}% denominators-of-projections-to-lifts
\chapters@set{20261126T031000}{40}% twists-with-few-congruences
\chapters@set{20261126T031100}{41}% theorem-c-and-the-main-inequality
\providecommand\chapnumber[1]{%
  \@ifundefined{chapters@#1}%
    {\PackageError{chapters}{Sheet #1 is not listed in the index}%
      {List it in 20261125T012000--grh-workbook__guide.org\MessageBreak
       and run scripts/make-chapters.py.}}%
    {\edef\thechap{\@nameuse{chapters@#1}}}}
\makeatother
\chapnumber{<sheet ID>}
\makeatletter
\providecommand\chapters@set[2]{\@namedef{chapters@#1}{#2}}
\chapters@set{20261125T001000}{1}% zeros-to-ramp-sums
\chapters@set{20261125T001100}{2}% inert-primes-and-estimate-b
\chapters@set{20261125T005000}{3}% galois-fourier-and-averaging
\chapters@set{20261125T005500}{4}% columns-from-a-sphere
\chapters@set{20261125T000000}{5}% slope-method
\chapters@set{20261126T030100}{6}% a-twist-for-each-prime
\chapters@set{20261125T003000}{7}% the-chain-and-the-square-root
\chapters@set{20261125T010000}{8}% rank-lemma-and-zero-estimate
\chapters@set{20261125T010500}{9}% poincare-series-and-near-orthogonality
\chapters@set{20261124T223000}{10}% elliptic-curves-for-the-proof
\chapters@set{20261125T009500}{11}% fock-model-and-theta-intertwiners
\chapters@set{20261124T223100}{12}% complex-multiplication-and-inert-reduction
\chapters@set{20261125T008000}{13}% one-cm-determinant
\chapters@set{20261124T214500}{14}% supersingular-theta-divisibility
\chapters@set{20261124T214600}{15}% the-isotropic-plane
\chapters@set{20261125T013700}{16}% pairings-on-a-formal-group-of-height-two
\chapters@set{20261125T013500}{17}% hecke-nondegeneracy-on-the-tree
\chapters@set{20261126T030200}{18}% the-bound-is-attained
\chapters@set{20261125T009000}{19}% kronecker-limit-and-chowla-selberg
\chapters@set{20261125T002000}{20}% definite-theta-lifts
\chapters@set{20261125T015500}{21}% siegel-weil-in-the-smallest-case
\chapters@set{20261125T002100}{22}% definite-theta-lifts-and-norms
\chapters@set{20261125T016500}{23}% hecke-operators-and-theta-lifts
\chapters@set{20261125T016700}{24}% divisibility-along-the-chain
\chapters@set{20261125T004000}{25}% denominators-and-sha
\chapters@set{20261125T004100}{26}% tate-twists-and-sha
\chapters@set{20261126T020430}{27}% the-minus-part-by-gauss-sums
\chapters@set{20261125T014100}{28}% congruence-modules
\chapters@set{20261125T014500}{29}% eisenstein-congruences-and-selmer-classes
\chapters@set{20261125T014700}{30}% pseudocharacters-and-the-length-bound
\chapters@set{20261125T011000}{31}% galois-representations-black-box
\chapters@set{20261126T030300}{32}% the-data
\chapters@set{20261126T020210}{33}% bessel-poincare-series-and-their-gram-operators
\chapters@set{20261126T030400}{34}% the-seesaw-identity
\chapters@set{20261126T030500}{35}% divisibility-of-the-composite-lift
\chapters@set{20261126T030600}{36}% theorem-a-short-forms
\chapters@set{20261126T030700}{37}% the-determinant-of-a-chain
\chapters@set{20261126T030800}{38}% the-cost-of-one-lift
\chapters@set{20261126T030900}{39}% denominators-of-projections-to-lifts
\chapters@set{20261126T031000}{40}% twists-with-few-congruences
\chapters@set{20261126T031100}{41}% theorem-c-and-the-main-inequality
\providecommand\chapnumber[1]{%
  \@ifundefined{chapters@#1}%
    {\PackageError{chapters}{Sheet #1 is not listed in the index}%
      {List it in 20261125T012000--grh-workbook__guide.org\MessageBreak
       and run scripts/make-chapters.py.}}%
    {\edef\thechap{\@nameuse{chapters@#1}}}}
\makeatother
\chapnumber{20261125T011000}
\renewcommand{\thesection}{\thechap.\arabic{section}}

\theoremstyle{definition}
\newtheorem{problem}{Problem}[section]
\newtheorem{definition}{Definition}[section]
\theoremstyle{plain}
\newtheorem{theorem}{Theorem}[section]

\newcommand{\Z}{\mathbb{Z}}
\newcommand{\Q}{\mathbb{Q}}
\newcommand{\R}{\mathbb{R}}
\newcommand{\C}{\mathbb{C}}
\newcommand{\F}{\mathbb{F}}
\newcommand{\A}{\mathbb{A}}
\newcommand{\Qlb}{\overline{\Q}_\ell}
\newcommand{\Oc}{\mathcal{O}}
\newcommand{\Gal}{\operatorname{Gal}}
\newcommand{\GL}{\mathrm{GL}}
\newcommand{\Art}{\operatorname{Art}}
\newcommand{\rec}{\operatorname{rec}}
\newcommand{\WD}{\operatorname{WD}}
\newcommand{\Frob}{\mathrm{Frob}}
\newcommand{\Ind}{\operatorname{Ind}}
\newcommand{\Sp}{\operatorname{Sp}}
\newcommand{\St}{\mathrm{St}}
\newcommand{\tr}{\operatorname{tr}}
\newcommand{\ii}{\mathrm{i}}

\setlength{\emergencystretch}{2em}
\title{Black boxes: Galois representations and local--global compatibility}
\author{}
\date{November 25, 2026}

\begin{document}
\maketitle
\refstepcounter{chap}\label{chap}

\paragraph{Prerequisites.}
Local and global class field theory, including ray class groups and the Chebotarev density theorem \cite[Ch.~V--VI, \S VII.13]{Ne}.  Elliptic curves at the level of \cite[Ch.~III, V]{Si1}; the curve \(y^2 = x^3 - x\) as in Chapter~\ref{cm-chap}, \emph{Complex multiplication and reduction at an inert prime}.  Modular forms of weight \(2\) and Hecke operators \cite[Ch.~1--6]{DS}.  Parabolic induction and the Steinberg representation of \(\GL_2\) over a \(p\)-adic field \cite[Ch.~3]{BH}.  The words crystalline, potentially semistable and Hodge--Tate weight are used only in statements.  Everything below that is not posed as a problem is a published theorem, stated with a reference, and is a black box for this chapter.

\section{Weil--Deligne representations and local Langlands maps}\label{sec:local}

What does a Galois representation remember at a prime \(p\ne\ell\), and what does an automorphic representation give to compare it with?

Let \(K/\Q_p\) be finite, with uniformizer \(\varpi_K\) and residue field of order \(q\).  Let \(W_K\) be its Weil group, \(I_K\) the inertia subgroup, and \(\Frob \in W_K\) a geometric Frobenius.  Let \(\Art_K \colon K^\times \to W_K^{ab}\) be the reciprocity map, normalized so that \(\varpi_K\) goes to a geometric Frobenius, and let \(\lvert w \rvert := \lvert \Art_K^{-1}(w) \rvert\), so that \(\lvert \Frob \rvert = q^{-1}\).  A \emph{Weil--Deligne representation} is a pair \((r, N)\) of a representation \(r \colon W_K \to \GL(V)\) on a finite-dimensional space over \(\C\) or \(\Qlb\), trivial on an open subgroup of \(I_K\), and an endomorphism \(N\) with \(r(w) N r(w)^{-1} = \lvert w \rvert N\).  It is \emph{Frobenius-semisimple} if \(r(\Frob)\) is semisimple.  Let \(\Sp(2)\) be the pair with \(V = \C^2\), \(r(w) = \mathrm{diag}(1, \lvert w \rvert)\) and \(N e_1 = e_2\), \(N e_2 = 0\).

\begin{problem}\label{pr:nilpotent}
In a Weil--Deligne representation, \(N\) is nilpotent.
\end{problem}

\begin{problem}\label{pr:sp2}
A two-dimensional Weil--Deligne representation with \(N \ne 0\) is isomorphic to \(\Sp(2) \otimes \chi\) for exactly one character \(\chi\) of \(W_K\).
\end{problem}

Let \(\ell \ne p\).  By Grothendieck's monodromy theorem \cite[4.2.1--4.2.2]{Ta}, every continuous \(\rho \colon W_K \to \GL_n(\Qlb)\) determines a Weil--Deligne representation \(\WD(\rho) = (r, N)\), characterized by \(\rho(\Frob^m \sigma) = r(\Frob^m \sigma) \exp(t_\ell(\sigma) N)\) for \(m \in \Z\) and \(\sigma \in I_K\), where \(t_\ell \colon I_K \to \Z_\ell\) is a surjection onto the \(\ell\)-part of tame inertia.  We write \(\WD(\rho)^{F\text{-ss}}\) for its Frobenius-semisimplification.

\begin{problem}\label{pr:monodromy}
If \(\rho(I_K)\) is finite, then \(N = 0\).  If \(\rho(I_K)\) is a nontrivial group of unipotent matrices, then \(N \ne 0\).
\end{problem}

For characters \(\chi_1, \dots, \chi_n\) of \(K^\times\), let \(i(\chi_1, \dots, \chi_n)\) be the normalized parabolic induction from the Borel subgroup of \(\GL_n(K)\).  Let \(\St\) be the Steinberg representation of \(\GL_2(K)\), the irreducible subquotient of \(i(\lvert\cdot\rvert^{1/2}, \lvert\cdot\rvert^{-1/2})\) that is not a character.

\begin{theorem}[Harris--Taylor, Henniart]\label{thm:LLC}
For each \(n \ge 1\) there is a bijection \(\rec_K\) from irreducible admissible representations of \(\GL_n(K)\) to Frobenius-semisimple \(n\)-dimensional Weil--Deligne representations of \(W_K\) such that \(\rec_K(\chi) = \chi \circ \Art_K^{-1}\) for \(n = 1\), \(\rec_K(\pi \otimes \chi \circ \det) = \rec_K(\pi) \otimes \rec_K(\chi)\), \(\det \rec_K(\pi) = \rec_K(\omega_\pi)\) for the central character \(\omega_\pi\), \(\rec_K(\pi^\vee) = \rec_K(\pi)^\vee\), and \(L\)- and \(\epsilon\)-factors of pairs are preserved.  These properties determine \(\rec_K\).
\end{theorem}

See \cite{HT}, \cite{He} and, for another proof, \cite{Sch}.  If \(i(\chi_1, \dots, \chi_n)\) is irreducible, then \(\rec_K\) sends it to \(\bigoplus_i \chi_i \circ \Art_K^{-1}\) with \(N = 0\); for every irreducible subquotient \(\pi\) of \(i(\chi_1, \dots, \chi_n)\), the semisimplification of the \(W_K\)-representation underlying \(\rec_K(\pi)\) is \(\bigoplus_i \chi_i \circ \Art_K^{-1}\); and \(\rec_K(\St \otimes \chi) = \Sp(2) \otimes \chi\lvert\cdot\rvert^{-1/2} \circ \Art_K^{-1}\) \cite{Ku}.  Only these properties are used below.

\section{Galois representations over CM fields}\label{sec:global}

Let \(F\) be a CM field with totally real subfield \(F^+\) and complex conjugation \(c\), let \(\ell\) be a prime and \(\iota \colon \Qlb \cong \C\).  Let \(\epsilon_\ell\) be the \(\ell\)-adic cyclotomic character.  A cuspidal automorphic representation \(\pi\) of \(\GL_a(\A_F)\) is \emph{polarized} if \(\pi^c \cong \pi^\vee \otimes (\chi \circ N_{F/F^+} \circ \det)\) for a character \(\chi\) of \(\A_{F^+}^\times/(F^+)^\times\) satisfying a sign condition at the real places, and \emph{regular algebraic} if \(\pi_\infty\) has the infinitesimal character of an irreducible algebraic representation of \(\mathrm{Res}_{F/\Q}\GL_a\); see \cite[\S2.1]{BLGGT}.

\begin{theorem}\label{thm:galois}
Let \(\pi\) be a regular algebraic polarized cuspidal automorphic representation of \(\GL_a(\A_F)\).  There is a continuous semisimple \(r = r_{\ell,\iota}(\pi) \colon G_F \to \GL_a(\Qlb)\) such that:
\textup{(i)} \(r^c \cong r^\vee \otimes \epsilon_\ell^{1-a} r_{\ell,\iota}(\chi)|_{G_F}\);
\textup{(ii)} for every finite place \(w \nmid \ell\), \(\iota\WD(r|_{G_{F_w}})^{F\text{-ss}} \cong \rec_{F_w}(\pi_w \otimes \lvert\det\rvert^{(1-a)/2})\);
\textup{(iii)} for every \(w \mid \ell\), \(r|_{G_{F_w}}\) is potentially semistable, crystalline if \(\pi_w\) is unramified, with Hodge--Tate weights determined by \(\pi_\infty\), and \textup{(ii)} holds for the Weil--Deligne representation attached to it by Fontaine.
\end{theorem}

The representations are constructed in the cohomology of Shimura varieties attached to unitary groups, to which \(\pi\) is first descended, by Kottwitz, Clozel, Harris--Taylor, Shin and Chenevier--Harris \cite{Cl,HT,Shi,CH}.  Part (ii) is due to Taylor--Yoshida and Caraiani \cite{TY,Car1}, and part (iii) to Barnet-Lamb--Gee--Geraghty--Taylor and Caraiani \cite[Thm.~2.1.1]{BLGGT}, \cite{Car2}.  For Alpöge's argument \cite{Alp}, the reconstruction \cite[\S17.5]{Adr} lists Theorem~\ref{thm:galois}, with local--global compatibility at all places, among the published inputs of its Galois-theoretic steps.

The smallest case is \(a = 1\) over an imaginary quadratic field \(K\), where the unitary group is \(U(1) := \ker(N \colon \mathrm{Res}_{K/\Q}\mathbb{G}_m \to \mathbb{G}_m)\).

\begin{problem}\label{pr:U1-local}
\(U(1)(\R)\) is the unit circle.  \(U(1)(\Q_p) \cong \Q_p^\times\) if \(p\) splits in \(K\), and \(U(1)(\Q_p)\) is compact otherwise.
\end{problem}

\begin{problem}\label{pr:U1-descent}
Let \(\phi\) be a character of \(\A_K^\times/K^\times\) that is trivial on \(\A_\Q^\times\).  Then \(\phi^c = \phi^{-1}\), and \(\phi = \mu(x/x^c)\) for exactly one character \(\mu\) of \(U(1)(\A)/U(1)(\Q)\).  Conversely each such \(\mu\) gives such a \(\phi\).
\end{problem}

\begin{problem}\label{pr:U1-sign}
If \(\phi^c = \phi^{-1}\), then \(\phi|_{\A_\Q^\times}\) is either trivial or the quadratic character \(\omega_{K/\Q}\).
\end{problem}

So the characters with \(\phi^c = \phi^{-1}\) fall into two classes, and by Problem~\ref{pr:U1-descent} only the first comes from \(U(1)\).

\section{The Hecke character of \texorpdfstring{\(\Q(\ii)\)}{Q(i)}}\label{sec:cm}

What is the representation of Theorem~\ref{thm:galois} when \(a = 1\), and how does it compare with the curve \(E \colon y^2 = x^3 - x\)?  Let \(K = \Q(\ii)\), \(\Oc = \Z[\ii]\) and \(\mathfrak{f} = (2 + 2\ii)\).

\begin{problem}\label{pr:primary}
Every ideal \(\mathfrak{a}\) of \(\Oc\) prime to \(1 + \ii\) has exactly one generator \(\alpha \equiv 1 \pmod{\mathfrak{f}}\).  Let \(\psi(\mathfrak{a}) := \alpha\).  Then \(\psi\) is multiplicative, \(\psi(\mathfrak{a}^c) = \overline{\psi(\mathfrak{a})}\), and \(\psi(\mathfrak{a})\psi(\mathfrak{a}^c) = N\mathfrak{a}\).
\end{problem}

For \(p \equiv 1 \pmod 4\) and \(\mathfrak{p} \mid p\), one has \(p + 1 - \#E(\F_p) = \psi(\mathfrak{p}) + \psi(\mathfrak{p}^c)\) \cite[Ch.~18, \S4]{IR}.

\begin{problem}\label{pr:psi-values}
\(\psi((1 + 2\ii)) = -1 - 2\ii\) and \(\psi((3 + 2\ii)) = 3 + 2\ii\), and \(\#E(\F_5) = \#E(\F_{13}) = 8\).
\end{problem}

\begin{problem}\label{pr:ray}
Let \(n \ge 1\).  If ideals \(\mathfrak{a}, \mathfrak{b}\) prime to \(2\ell\) have the same class in the ray class group modulo \(\mathfrak{f}\ell^n\), then \(\psi(\mathfrak{a}) \equiv \psi(\mathfrak{b}) \pmod{\ell^n\Oc}\).
\end{problem}

Combining Problem~\ref{pr:ray} with Artin reciprocity for the ray class fields modulo \(\mathfrak{f}\ell^n\) gives the following.

\begin{theorem}\label{thm:psi-l}
There is a continuous character \(\psi_\ell \colon G_K \to (\Oc \otimes \Z_\ell)^\times\), unramified outside \(2\ell\), with \(\psi_\ell(\Frob_\mathfrak{p}) = \psi(\mathfrak{p})\) for every prime \(\mathfrak{p} \nmid 2\ell\), where \(\Frob_\mathfrak{p}\) is an arithmetic Frobenius.
\end{theorem}

By Problem~\ref{pr:primary}, \(\psi\psi^c\) is the norm, so \(\psi\) is polarized, and \(\psi_\ell\), followed by an embedding \(\Oc \otimes \Z_\ell \to \Qlb\), is the representation of Theorem~\ref{thm:galois} for \(a = 1\), up to the normalizations of the reciprocity maps.  At the places not above \(\ell\) the theorem is then class field theory; the next problem looks at the places above \(\ell\).

\begin{problem}\label{pr:HT}
Let \(\ell \equiv 1 \pmod 4\), \(\ell = \lambda\bar\lambda\), and let \(\psi_\lambda \colon G_K \to \Z_\ell^\times\) be \(\psi_\ell\) followed by \(\Oc \otimes \Z_\ell \to \Oc_\lambda = \Z_\ell\).  Then \(\psi_\lambda = \epsilon_\ell\) on an open subgroup of the inertia group at \(\lambda\), and \(\psi_\lambda\) has finite image on the inertia group at \(\bar\lambda\).
\end{problem}

So \(\psi_\lambda\) has Hodge--Tate weights \(1\) and \(0\) at the two places above \(\ell\) (in the convention where \(\epsilon_\ell\) has weight \(1\)), and the polarization appears as \(\mathrm{HT}_{\tau c} = 1 - \mathrm{HT}_\tau\).

The automorphism \([\ii](x, y) = (-x, \ii y)\) of \(E\) is defined over \(K\), so \(V_\ell(E) := T_\ell(E) \otimes \Q\) is free of rank one over \(K \otimes \Q_\ell\), and \(G_K\) acts on it through a character \(G_K \to (K \otimes \Q_\ell)^\times\).

\begin{problem}\label{pr:frob-cm}
Let \(\mathfrak{p} \nmid 2\ell\).  If \(N\mathfrak{p} = p \equiv 1 \pmod 4\), then \(\Frob_\mathfrak{p}\) acts on \(V_\ell(E)\) as \(\psi(\mathfrak{p})\).  If \(\mathfrak{p} = (p)\) with \(p \equiv 3 \pmod 4\), it acts as \(-p = \psi((p))\).
\end{problem}

By the Chebotarev density theorem, Problem~\ref{pr:frob-cm} proves the following; it is due to Deuring, see \cite[Ch.~II, \S10]{Si2}.

\begin{theorem}\label{thm:deuring}
The character of \(G_K\) on \(V_\ell(E)\) is \(\psi_\ell\).
\end{theorem}

\begin{problem}\label{pr:induced}
Let \(\psi_1, \psi_2 \colon G_K \to \Qlb^\times\) be \(\psi_\ell\) followed by the two embeddings \(\Oc \otimes \Z_\ell \to \Qlb\).  Then \(\psi_2 = \psi_1^c\), \(V_\ell(E) \otimes \Qlb \cong \Ind_{G_K}^{G_\Q} \psi_1\), and \(\tr(\Frob_p \mid V_\ell(E)) = 0\) for \(p \equiv 3 \pmod 4\), \(p \ne \ell\).
\end{problem}

Restricted to \(G_K\), \(V_\ell(E)\) is \(\psi_1 \oplus \psi_1^c\), which is reducible.  So the weight-two form of \(E\), base changed to \(K\), is not cuspidal, and it enters Theorem~\ref{thm:galois} only through \(\psi\), with \(a = 1\).  For a cuspidal example with \(a = 2\) one needs a curve without complex multiplication by \(K\).  We take \(X_0(11)\) and stay over \(\Q\); its base change to \(K\) is cuspidal, and the local questions are the same.

\section{Eichler--Shimura for \texorpdfstring{\(X_0(11)\)}{X0(11)}}\label{sec:es}

How does one get Frobenius eigenvalues out of the cohomology of a Shimura variety, and what does each method give?  Let
\[
  f = q \prod_{n \ge 1} (1 - q^n)^2 (1 - q^{11n})^2 = \sum_{n \ge 1} a_n q^n ,
\]
the normalized eigenform spanning \(S_2(\Gamma_0(11))\) \cite[\S3.9]{DS}.  The curve \(X_0(11)\) has genus one and the model \(E \colon y^2 + y = x^3 - x^2 - 10x - 20\) over \(\Q\), with good reduction away from \(11\).  For \(p \ne 11\), the Hecke operator \(T_p\) acts on \(E = J_0(11)\) as \([a_p]\), and by the Eichler--Shimura relation \cite[Thm.~8.7.2]{DS}, \cite[Thm.~7.9]{Sh}, the reduction of \(T_p\) modulo \(p\) is \(F + V\), where \(F\) is the Frobenius endomorphism and \(FV = VF = p\).

\begin{problem}\label{pr:coeffs}
\(a_2 = -2\), \(a_3 = -1\), \(a_5 = 1\), \(a_7 = -2\), \(a_{11} = 1\), \(a_{13} = 4\).
\end{problem}

\begin{problem}\label{pr:annihilate}
For \(p \notin \{11, \ell\}\), the arithmetic Frobenius at \(p\) on \(V_\ell(E)\) satisfies \(\Frob_p^2 - a_p\Frob_p + p = 0\).
\end{problem}

For \(p \notin \{11, \ell\}\), the arithmetic Frobenius at \(p\) has determinant \(p\) on \(V_\ell(E)\), by the Weil pairing \cite[\S III.8]{Si1}, and \(\#E(\F_{p^k}) = p^k + 1 - \tr(\Frob_p^k \mid V_\ell(E))\) for \(k \ge 1\) \cite[\S V.2]{Si1}.

\begin{problem}\label{pr:charpoly}
Let \(\Phi \in \GL_2(\Qlb)\) satisfy \(\Phi^2 - a\Phi + p = 0\) with \(a \in \Z\), and let \(\det\Phi = p\).  Then the characteristic polynomial of \(\Phi\) is \(X^2 - aX + p\).  Hence, taking \(\ell \ne p\), \(\#E(\F_p) = p + 1 - a_p\) for \(p \ne 11\).
\end{problem}

The relation of Problem~\ref{pr:annihilate} kills Frobenius; it does not give the multiplicities of its roots.  In dimension two the determinant supplies them.  In dimension three, \(\mathrm{diag}(\alpha, \alpha, \beta)\) and \(\mathrm{diag}(\alpha, \beta, \beta)\) are killed by the same quadratic polynomial.  The published constructions behind Theorem~\ref{thm:galois} obtain the whole characteristic polynomial by counting points, following Langlands and Kottwitz \cite{Ko}.  The simplest instance is the trace formula for \(\#E(\F_{p^k})\) stated before Problem~\ref{pr:charpoly}.

\begin{problem}\label{pr:count}
\(\#E(\F_p) = p + 1 - a_p\) for \(p = 2, 3, 5, 7, 13\) by direct count.  The numbers \(\#E(\F_p)\) and \(\#E(\F_{p^2})\) determine the characteristic polynomial of \(\Frob_p\); for \(p = 2\) both equal \(5\).
\end{problem}

At \(11\) the curve has split multiplicative reduction.  By Tate's theory \cite[Ch.~V]{Si2}, \(E \cong \mathbb{G}_m/q^\Z\) over \(\Q_{11}\) for some \(q \in 11\Z_{11}\), and \(T_\ell(E)\) is an extension of \(\Z_\ell\) by \(\Z_\ell(1)\) on which the inertia group at \(11\) acts through the Kummer classes of \(q^{1/\ell^n}\).

\begin{problem}\label{pr:tate}
\(\WD(V_\ell(E)|_{G_{\Q_{11}}}) \cong \Sp(2) \otimes \chi\) for an unramified character \(\chi\); in particular \(N \ne 0\).
\end{problem}

Correspondingly, \(\pi_{11}\) is an unramified twist of \(\St\), and the compatibility of Theorem~\ref{thm:galois}(ii), in its version over \(\Q\) due to Carayol \cite{Ca}, matches the monodromy at \(11\).  The later parts, from Chapter~\ref{ob-chap} on, state each ingredient in the generality the proof uses and assemble them.

\section*{Hints}

\subsection*{\S\ref{sec:local}}
\begin{description}
\item[\ref{pr:nilpotent}] Take \(w = \Frob\): \(N\) is conjugate to \(q^{-1}N\).
\item[\ref{pr:sp2}] \(\ker N\) is a line stable under \(W_K\); compare the characters on \(\ker N\) and on \(V/\ker N\).  The extension splits because \(\lvert\Frob\rvert \ne 1\).
\item[\ref{pr:monodromy}] \(\exp(t_\ell(\sigma)N) = 1\) on an open subgroup of \(I_K\) exactly when \(N = 0\).
\end{description}

\subsection*{\S\ref{sec:global}}
\begin{description}
\item[\ref{pr:U1-local}] \(K \otimes \Q_p\) is \(\Q_p \times \Q_p\) or a field.
\item[\ref{pr:U1-descent}] Hilbert's Theorem 90 for \(K\) and for each \(K \otimes \Q_v\), and for the units at the unramified places. % check-numbers: ok (Hilbert): Theorem 90
\item[\ref{pr:U1-sign}] \(\phi\) is trivial on \(\Q^\times\) and on the norms \(x x^c\); use \(\A_\Q^\times/\Q^\times N(\A_K^\times) \cong \Z/2\).
\end{description}

\subsection*{\S\ref{sec:cm}}
\begin{description}
\item[\ref{pr:primary}] The units \(\pm 1, \pm\ii\) map bijectively onto \((\Oc/\mathfrak{f})^\times\), which has order \(4\); and \(\mathfrak{f}^c = \mathfrak{f}\).
\item[\ref{pr:ray}] Write \(\mathfrak{a} = \mathfrak{b}(\beta)\) with \(\beta \equiv 1 \pmod{^\times \mathfrak{f}\ell^n}\); then \(\beta\) is the generator of \((\beta)\) that \(\psi\) picks.
\item[\ref{pr:HT}] Approximate an idèle \(u \in 1 + \ell^m\Z_\ell\) placed at \(\lambda\) by a global element of \(\Oc\).  With uniformizers sent to arithmetic Frobenius, as in Theorem~\ref{thm:psi-l}, the local reciprocity map sends \(u\) to an element acting on \(\ell\)-power roots of unity by \(u^{-1}\).
\item[\ref{pr:frob-cm}] The \(p\)-power Frobenius of the reduction commutes with \([\ii]\), so it is \([\alpha]\) with \(\alpha \in \Oc\) of norm \(p\).  It is inseparable, so \(\alpha \in \mathfrak{p}\); fix the unit with \(a_p = \psi(\mathfrak{p}) + \psi(\mathfrak{p}^c)\).  For \(p \equiv 3 \pmod 4\), use \(a_p = 0\).
\item[\ref{pr:induced}] Complex conjugation acts on \(V_\ell(E)\) and anticommutes with \([\ii]\).
\end{description}

\subsection*{\S\ref{sec:es}}
\begin{description}
\item[\ref{pr:coeffs}] Expand the product to order \(q^{13}\).
\item[\ref{pr:annihilate}] Multiply \(T_p = F + V\) by \(F\), and use that \(\Frob_p\) acts on \(T_\ell\) of the reduction as \(F\).
\item[\ref{pr:charpoly}] Treat scalar \(\Phi\) separately.
\item[\ref{pr:count}] In \(\F_4\), the equation becomes \(y^2 + y = x^3 + x^2\).  Use \(\alpha^2 + \beta^2 = (\alpha+\beta)^2 - 2\alpha\beta\).
\item[\ref{pr:tate}] Problems~\ref{pr:sp2} and~\ref{pr:monodromy}; \(v_{11}(q) \ne 0\).
\end{description}

\begin{thebibliography}{BLGGT}
\bibitem[Adr]{Adr} J.~Adriazola, \emph{The Alp\"oge proof of the Riemann hypothesis}, manuscript, revised 24 November 2026, \url{https://github.com/jimjamzola/RH-Alpoge-Reconstruction}.
\bibitem[Alp]{Alp} L.~Alp\"oge, untitled note announcing a proof of the Riemann hypothesis, 20 November 2026, \url{https://alpo.ge/remark.pdf}.
\bibitem[BH]{BH} C.~J.~Bushnell and G.~Henniart, \emph{The Local Langlands Conjecture for \(\mathrm{GL}(2)\)}, Grundlehren Math.\ Wiss.\ 335, Springer, 2006.
\bibitem[BLGGT]{BLGGT} T.~Barnet-Lamb, T.~Gee, D.~Geraghty and R.~Taylor, \emph{Potential automorphy and change of weight}, Ann.\ of Math.\ (2) \textbf{179} (2014), 501--609.
\bibitem[Ca]{Ca} H.~Carayol, \emph{Sur les repr\'esentations \(\ell\)-adiques associ\'ees aux formes modulaires de Hilbert}, Ann.\ Sci.\ \'Ec.\ Norm.\ Sup.\ (4) \textbf{19} (1986), 409--468.
\bibitem[Car1]{Car1} A.~Caraiani, \emph{Local--global compatibility and the action of monodromy on nearby cycles}, Duke Math.\ J.\ \textbf{161} (2012), 2311--2413.
\bibitem[Car2]{Car2} A.~Caraiani, \emph{Monodromy and local--global compatibility for \(l = p\)}, Algebra Number Theory \textbf{8} (2014), 1597--1646.
\bibitem[CH]{CH} G.~Chenevier and M.~Harris, \emph{Construction of automorphic Galois representations, II}, Camb.\ J.\ Math.\ \textbf{1} (2013), 53--73.
\bibitem[Cl]{Cl} L.~Clozel, \emph{Repr\'esentations galoisiennes associ\'ees aux repr\'esentations automorphes autoduales de \(\mathrm{GL}(n)\)}, Publ.\ Math.\ IH\'ES \textbf{73} (1991), 97--145.
\bibitem[DS]{DS} F.~Diamond and J.~Shurman, \emph{A First Course in Modular Forms}, Grad.\ Texts in Math.\ 228, Springer, 2005.
\bibitem[He]{He} G.~Henniart, \emph{Une preuve simple des conjectures de Langlands pour \(\mathrm{GL}(n)\) sur un corps \(p\)-adique}, Invent.\ Math.\ \textbf{139} (2000), 439--455.
\bibitem[HT]{HT} M.~Harris and R.~Taylor, \emph{The Geometry and Cohomology of Some Simple Shimura Varieties}, Ann.\ of Math.\ Stud.\ 151, Princeton Univ.\ Press, 2001.
\bibitem[IR]{IR} K.~Ireland and M.~Rosen, \emph{A Classical Introduction to Modern Number Theory}, 2nd ed., Grad.\ Texts in Math.\ 84, Springer, 1990.
\bibitem[Ko]{Ko} R.~E.~Kottwitz, \emph{Points on some Shimura varieties over finite fields}, J.\ Amer.\ Math.\ Soc.\ \textbf{5} (1992), 373--444.
\bibitem[Ku]{Ku} S.~S.~Kudla, \emph{The local Langlands correspondence: the non-archimedean case}, in: Motives, Proc.\ Sympos.\ Pure Math.\ 55, Part~2, Amer.\ Math.\ Soc., 1994, 365--391.
\bibitem[Ne]{Ne} J.~Neukirch, \emph{Algebraic Number Theory}, Grundlehren Math.\ Wiss.\ 322, Springer, 1999.
\bibitem[Sch]{Sch} P.~Scholze, \emph{The local Langlands correspondence for \(\mathrm{GL}_n\) over \(p\)-adic fields}, Invent.\ Math.\ \textbf{192} (2013), 663--715.
\bibitem[Sh]{Sh} G.~Shimura, \emph{Introduction to the Arithmetic Theory of Automorphic Functions}, Princeton Univ.\ Press, 1971.
\bibitem[Shi]{Shi} S.~W.~Shin, \emph{Galois representations arising from some compact Shimura varieties}, Ann.\ of Math.\ (2) \textbf{173} (2011), 1645--1741.
\bibitem[Si1]{Si1} J.~H.~Silverman, \emph{The Arithmetic of Elliptic Curves}, 2nd ed., Grad.\ Texts in Math.\ 106, Springer, 2009.
\bibitem[Si2]{Si2} J.~H.~Silverman, \emph{Advanced Topics in the Arithmetic of Elliptic Curves}, Grad.\ Texts in Math.\ 151, Springer, 1994.
\bibitem[Ta]{Ta} J.~Tate, \emph{Number theoretic background}, in: Automorphic Forms, Representations and \(L\)-Functions, Proc.\ Sympos.\ Pure Math.\ 33, Part~2, Amer.\ Math.\ Soc., 1979, 3--26.
\bibitem[TY]{TY} R.~Taylor and T.~Yoshida, \emph{Compatibility of local and global Langlands correspondences}, J.\ Amer.\ Math.\ Soc.\ \textbf{20} (2007), 467--493.
\end{thebibliography}

\end{document}
